\documentclass[../../main.tex]{subfiles}
\begin{document}

{\bf[解]} \1 $O$を原点とし，水平線を$x$軸とする下図のような座標系で考える．

     \begin{minipage}{0.5\hsize}
          \begin{center}
          \scalebox{.4}{\input{ut-62-3p1}}
          \end{center}
     \end{minipage}
     \begin{minipage}{0.5\hsize}
         \begin{center}
          \scalebox{.5}{\input{ut-62-3p2}}
          \end{center}
     \end{minipage}
     
$OA=1$として考えて一般性を失わない．また$G$が最低のときを考えるので，$A$は$g$の下側にあるとしてよい．
簡単のため$\b=\theta-\a$とおく．$A$から$g$に下ろした垂足$H$とすると，$OA=AB$ゆえ，$G$は$AH$の中点である．
$OH=\cos\b$ゆえ，
     \begin{align}
     &A(\cos(-\theta),\sin(-\theta)) \\
     &H(\cos\b\cos(-\a),\cos\b\sin(-\b))
     \end{align}
となるから，
     \begin{align}
     \bekutoru{OG}&=\bekutoru{OH}+\dfrac{1}{2}\bekutoru{HA} \\
     &=\dfrac{1}{2}\bekutoru{OH}+\dfrac{1}{2}\bekutoru{OA} \\
     &=\dfrac{1}{2}\cos\b\vtwo{\cos\a}{-\sin\a}+\dfrac{1}{2}\vtwo{c}{-s}
     \end{align}
である．この$y$座標$Y$として
     \begin{align}
     Y=\dfrac{-1}{2}(\sin\a\cos\b+s)
     \end{align}
だから，これが最小になる時の$\tan\theta$を求めればよい．
     \[\a+\b=\theta\]
に注意して
     \begin{align}
     -2Y&=\dfrac{1}{2}(\sin(\a+\b)+\sin(\a-\b))+s \\
     &=\dfrac{3}{2}s+\dfrac{1}{2}\sin(2\a-\theta)\atag\label{1}
     \end{align}
題意から$\tan\a=1/2$であるから
     \begin{align}
     \sin2\a=\dfrac{2\tan\a}{1+\tan^2\a}=\dfrac{4}{5} \\
     \cos2\a=\dfrac{1-\tan^2\a}{1+\tan^2\a}=\dfrac{3}{5}
     \end{align}
が従う．\eqref{1}に代入して
     \begin{align}
     Y&=\dfrac{-1}{4}[3s+c\sin2\a -s\cos2\a ] \\
     &=\dfrac{-1}{4}[3s+\dfrac{4}{5} c-\dfrac{3}{5} s] \\
     &=\dfrac{-1}{5}(3s+c) \\
     &\ge \dfrac{-1}{5}\sqrt{3^2+1}&(\because \text{コーシー})
     \end{align}
統合成立条件は
     \begin{align}
     &\vtwo{c}{s}\parallel\vtwo{1}{3} \\
     &s-3c=0
     \end{align}
の時で，この時
     \[\tan\theta=3\]
である．$\cdots$(答)
     
     
\newpage
\end{document}